\chapter{Vector Spaces and Subspaces}

So far we have been doing linear algebra with our sleeves rolled up. We have added
vectors in $\R^2$ and $\R^3$, scaled them, taken their dot products, packed them into
matrices, and pushed those matrices through elimination. Every object we touched was a
list of numbers. That was the right place to start, because lists of numbers are
something you can \emph{see}: an arrow on a page, a point in space, a recipe for mixing
ingredients.

Now we take a step back and ask a quieter question. Forget for a moment
\emph{what} a vector is made of. Ask instead: what can we \emph{do} with vectors, and
which collections of them are ``closed'' under those operations? The answer is the idea
of a \emph{vector space}, and its most useful relative, the \emph{subspace}.

This is the chapter where linear algebra stops being a bag of tricks for $\R^n$ and
becomes a language. Once you have the vector-space idea, the same theorems apply to
arrows, to functions, to signals, to polynomials, to images, to the solution sets of
differential equations. You will prove something once and use it everywhere. We will
build the idea slowly, picture first.

\chapterline

\section{What ``space'' is supposed to capture}

\subsection{The two operations we keep using}

Look back at everything we did with vectors. Strip it down and only two operations ever
appeared:

\begin{itemize}
  \item \textbf{Adding} two vectors to get a third vector: $\vec{u} + \vec{v}$.
  \item \textbf{Scaling} a vector by a number: $c\,\vec{v}$.
\end{itemize}

Dot products, lengths, and angles are wonderful, but they are \emph{extra}. The bare
skeleton of linear algebra is just ``add'' and ``scale.'' A linear combination
$c_1\vec{v}_1 + c_2\vec{v}_2 + \dots + c_k\vec{v}_k$ is nothing but these two operations
used together. Span, independence, dimension, rank, and eigenvectors will all be defined
using only adding and scaling.

\begin{intuitionbox}
A vector space is any world where ``add'' and ``scale'' both make sense and never throw
you out of the world. You start inside the world, you add or scale, and you always land
back inside. Nothing leaks out. That ``nothing leaks out'' property -- called
\emph{closure} -- is the whole game.
\end{intuitionbox}

The numbers we scale by are called \emph{scalars}. For us they will always be real
numbers (we will say a word about complex scalars later in the book). A vector space whose
scalars are real numbers is called a \emph{real vector space}.

\subsection{Why bother abstracting?}

Here is the payoff stated up front, before we have even seen the definition. Consider
three collections that look nothing alike:

\begin{itemize}
  \item all arrows in the plane $\R^2$;
  \item all polynomials of degree at most $2$, like $3 - x + 5x^2$;
  \item all smooth functions $y(t)$ satisfying $y'' + y = 0$.
\end{itemize}

The first is geometry, the second is algebra, the third is calculus. Yet in each one you
can add two members and get a member, and scale a member and get a member. They are all
vector spaces. Anything we prove about ``vector spaces in general'' is automatically true
in all three at once. That is the leverage. You do the work once.

\begin{aestheticbox}
There is a deep pleasure in the moment you realize that ``a function'' and ``an arrow''
can be the \emph{same kind of thing}. A sound wave, a grayscale image, a probability
distribution, a quantum state -- all of them live in vector spaces, and all of them obey
the theorems you are about to meet. Linear algebra is the mathematics of everything that
adds and scales, which turns out to be a startling fraction of everything.
\end{aestheticbox}

\section{The axioms of a vector space (made friendly)}

We now write down the official rules. Do not be intimidated by the list; every line is
something you already believe about ordinary arrows. The point of writing them out is so
that later we can \emph{check} whether some new, exotic collection deserves to be called
a vector space.

\begin{formalbox}
A \textbf{(real) vector space} is a set $V$ together with two operations -- vector
addition $\vec{u} + \vec{v}$ and scalar multiplication $c\,\vec{v}$ for $c \in \R$ --
such that for all $\vec{u},\vec{v},\vec{w} \in V$ and all scalars $c,d \in \R$:

\smallskip
\textbf{Closure.}
\begin{enumerate}[label=(\arabic*)]
  \item $\vec{u} + \vec{v} \in V$ \hfill (sums stay inside)
  \item $c\,\vec{v} \in V$ \hfill (scalings stay inside)
\end{enumerate}

\textbf{Addition behaves.}
\begin{enumerate}[label=(\arabic*),start=3]
  \item $\vec{u} + \vec{v} = \vec{v} + \vec{u}$ \hfill (commutative)
  \item $(\vec{u} + \vec{v}) + \vec{w} = \vec{u} + (\vec{v} + \vec{w})$ \hfill (associative)
  \item there is a \emph{zero vector} $\vec{0}\in V$ with $\vec{v} + \vec{0} = \vec{v}$
  \item every $\vec{v}$ has a \emph{negative} $-\vec{v}\in V$ with $\vec{v} + (-\vec{v}) = \vec{0}$
\end{enumerate}

\textbf{Scaling behaves.}
\begin{enumerate}[label=(\arabic*),start=7]
  \item $1\,\vec{v} = \vec{v}$ \hfill (scaling by one does nothing)
  \item $c(d\,\vec{v}) = (cd)\,\vec{v}$ \hfill (scalings compose)
  \item $c(\vec{u} + \vec{v}) = c\,\vec{u} + c\,\vec{v}$ \hfill (distribute over vectors)
  \item $(c + d)\,\vec{v} = c\,\vec{v} + d\,\vec{v}$ \hfill (distribute over scalars)
\end{enumerate}
\end{formalbox}

That is ten rules. Eight of them (3 through 10) are just ``the algebra you already do
works the way you expect.'' The only two with real teeth are the first two: closure under
addition and closure under scaling. Keep your eye on those.

\begin{focusbox}
When you meet a new candidate vector space, the rules that almost always do the work are
\textbf{closure under addition} and \textbf{closure under scalar multiplication}, plus
\textbf{``is $\vec{0}$ actually in there?''} The eight bookkeeping rules are usually
inherited for free from the arithmetic of real numbers. Spend your energy on closure.
\end{focusbox}

\subsection{Two small theorems you get for free}

The axioms are spare on purpose. Notice they do \emph{not} say ``$0\,\vec{v} = \vec{0}$''
or ``$(-1)\vec{v} = -\vec{v}$.'' Those feel obvious, but in a properly austere theory we
should \emph{derive} them. Watch how the rules force them.

\begin{proposition}\label{ch11:zero-scalar}
In any vector space, $0\,\vec{v} = \vec{0}$ for every $\vec{v}$.
\end{proposition}

\begin{proof}
Start from a true fact about real numbers, $0 + 0 = 0$, and hit $\vec{v}$ with it. By
rule (10),
\[
0\,\vec{v} = (0+0)\,\vec{v} = 0\,\vec{v} + 0\,\vec{v}.
\]
Now add the negative of $0\,\vec{v}$ (which exists by rule (6)) to both sides. The left
becomes $\vec{0}$, and on the right $0\,\vec{v} + 0\,\vec{v} + (-(0\,\vec{v}))$ collapses
to $0\,\vec{v}$. So $\vec{0} = 0\,\vec{v}$. Done.
\end{proof}

\begin{proposition}\label{ch11:neg-one}
In any vector space, $(-1)\,\vec{v} = -\vec{v}$.
\end{proposition}

\begin{proof}
We show $(-1)\vec{v}$ is \emph{the} negative of $\vec{v}$, i.e. it adds to $\vec{0}$:
\[
\vec{v} + (-1)\vec{v} = 1\,\vec{v} + (-1)\vec{v} = (1 + (-1))\,\vec{v} = 0\,\vec{v} = \vec{0},
\]
using rule (7), then rule (10), then arithmetic, then
Proposition~\ref{ch11:zero-scalar}. Since negatives are unique (a one-line exercise),
$(-1)\vec{v} = -\vec{v}$.
\end{proof}

\begin{intuitionbox}
These tiny proofs are your first taste of the abstract style. We never used coordinates,
never wrote a single number inside a vector. We used only the rules. That is exactly why
the conclusion holds in \emph{every} vector space at once -- arrows, polynomials,
functions, all of them. Coordinates are a convenience; the rules are the law.
\end{intuitionbox}

\section{The home example: $\R^n$}

The first vector space, and the one you should always picture, is $\R^n$: all columns of
$n$ real numbers,
\[
\R^n = \left\{ \colvec{x_1\\x_2\\\vdots\\x_n} : x_1,\dots,x_n \in \R \right\},
\]
with the addition and scaling you already know -- componentwise. For $n=2$ and $n=3$ you
can literally draw the vectors as arrows from the origin. For larger $n$ you cannot draw
them, but the rules are identical, so your geometric intuition keeps working.

\begin{example}\label{ch11:r2-check}
Let us sanity-check that $\R^2$ obeys closure. Take
$\vec{u} = \colvec{2\\-1}$ and $\vec{v} = \colvec{-3\\4}$ and the scalar $c = 5$. Then
\[
\vec{u} + \vec{v} = \colvec{2 + (-3)\\ -1 + 4} = \colvec{-1\\3} \in \R^2,
\qquad
5\,\vec{u} = \colvec{10\\-5} \in \R^2.
\]
A sum of two columns of two reals is a column of two reals; a scalar times such a column
is again such a column. Closure holds, and you can see the eight bookkeeping rules are
just facts about adding and multiplying the entries.
\end{example}

\begin{figure}[H]
\centering
\begin{tikzpicture}[scale=0.95,>=Latex]
  \draw[help lines,gray!30] (-3.5,-1.5) grid (3.5,4.5);
  \draw[->] (-3.7,0)--(3.7,0) node[right]{$x_1$};
  \draw[->] (0,-1.7)--(0,4.7) node[above]{$x_2$};
  % u
  \draw[->,very thick,vecblue] (0,0)--(2,-1) node[below right]{$\vec{u}$};
  % v
  \draw[->,very thick,vecred] (0,0)--(-3,4) node[above left]{$\vec{v}$};
  % u+v
  \draw[->,very thick,structurepurple] (0,0)--(-1,3) node[above]{$\vec{u}+\vec{v}$};
  % parallelogram guides
  \draw[dashed,gray] (2,-1)--(-1,3);
  \draw[dashed,gray] (-3,4)--(-1,3);
\end{tikzpicture}
\caption{Addition in $\R^2$ is the parallelogram rule: slide $\vec{v}$ to the tip of
$\vec{u}$ and the diagonal is $\vec{u}+\vec{v}$. The sum is another arrow from the origin
-- still inside $\R^2$. That is closure, drawn.}
\label{ch11:fig-r2add}
\end{figure}

\begin{strangbox}
Strang opens his course with $\R^n$ and the four fundamental subspaces, and only later --
sometimes in a footnote -- mentions the general axioms. Reading this \emph{before} Strang,
take the axioms as the ``why $\R^n$ works'' story so that when he says ``column space'' or
``nullspace'' you already hear ``a subspace, i.e. a smaller vector space living inside
$\R^n$.'' Reading this \emph{after} Strang, the axioms are what let you carry every result
you learned about $\R^n$ over to function spaces and operator theory without redoing a
thing.
\end{strangbox}

\section{Spaces that are not columns of numbers}

To feel the power of the definition, we need examples where the ``vectors'' are not
columns at all. These are not curiosities; they are where modern applications live.

\subsection{Matrices form a vector space}

The set of all $2\times 2$ real matrices, call it $M_{2\times 2}$, is a vector space. You
add matrices entrywise and scale them entrywise, and both stay $2\times 2$. The zero
vector is the zero matrix $\begin{bmatrix}0&0\\0&0\end{bmatrix}$.

\begin{example}\label{ch11:matrix-space}
With $\mat{A} = \begin{bmatrix}1&2\\0&-1\end{bmatrix}$,
$\mat{B} = \begin{bmatrix}3&0\\1&4\end{bmatrix}$, and $c = 2$:
\[
\mat{A}+\mat{B} = \begin{bmatrix}4&2\\1&3\end{bmatrix},
\qquad
2\mat{A} = \begin{bmatrix}2&4\\0&-2\end{bmatrix}.
\]
Both are again $2\times 2$ matrices. Here a single ``vector'' is a whole $2\times 2$ grid
of numbers -- four numbers, so this space behaves a lot like $\R^4$ wearing a square
costume.
\end{example}

\subsection{Polynomials form a vector space}

Let $P_2$ be all polynomials of degree at most $2$:
\[
p(x) = a_0 + a_1 x + a_2 x^2, \qquad a_0,a_1,a_2 \in \R.
\]
Add two such polynomials and you get another one of degree at most $2$; scale one and the
degree cannot grow. The zero vector is the constant polynomial $0$. Closure holds.

\begin{example}\label{ch11:poly-space}
With $p = 3 - x + 5x^2$ and $q = 1 + 2x - 5x^2$:
\[
p + q = 4 + x + 0\,x^2 = 4 + x.
\]
Notice the $x^2$ terms cancelled and the sum has degree $1$ -- still ``degree at most
$2$,'' so still inside $P_2$. (This is exactly why we say ``at most.'' The polynomials of
degree \emph{exactly} $2$ are \emph{not} a vector space, because two of them can sum to
something of lower degree, and worse, the zero polynomial is missing.)
\end{example}

\subsection{Functions form a vector space}

This is the example that opens the most doors. Let $C[0,1]$ be all continuous functions $f$ on the interval
$[0,1]$. Add two continuous functions pointwise, $(f+g)(x) = f(x)+g(x)$, and the result is
continuous; scale a continuous function and it stays continuous. The zero vector is the
function that is $0$ everywhere. So $C[0,1]$ is a vector space -- an infinite-dimensional
one, but a vector space all the same.

\begin{figure}[H]
\centering
\begin{tikzpicture}[scale=1.0,>=Latex]
  \begin{axis}[
      width=11cm,height=5.4cm,
      axis lines=middle,
      xlabel={$x$}, ylabel={},
      xmin=0,xmax=1,ymin=-0.2,ymax=1.6,
      xtick={0,1},ytick=\empty,
      domain=0:1,samples=80,
      legend pos=outer north east,
      legend style={draw=none}]
    \addplot[vecblue,very thick] {x};            \addlegendentry{$f(x)=x$}
    \addplot[vecred,very thick] {sin(deg(pi*x))*0.5}; \addlegendentry{$g(x)=\tfrac12\sin(\pi x)$}
    \addplot[structurepurple,very thick,dashed] {x + sin(deg(pi*x))*0.5}; \addlegendentry{$f+g$}
  \end{axis}
\end{tikzpicture}
\caption{Adding functions is adding arrows ``one $x$ at a time.'' At each input $x$ you
add the two output heights. The dashed curve $f+g$ is the sum vector. Continuous plus
continuous is continuous, so we never leave the space.}
\label{ch11:fig-funcadd}
\end{figure}

\begin{applicationbox}
\textbf{Where function spaces run the modern world.} A grayscale image is a function from
pixels to brightness; adding two images (double exposure) and scaling brightness are
exactly the vector-space operations, which is why a $64\times 64$ image lives in
$\R^{4096}$ and can be fed to a neural network as a vector. An audio clip is a function
from time to amplitude; mixing two tracks is vector addition, turning the volume knob is
scalar multiplication. The Fourier transform, JPEG compression, and the embeddings inside
a large language model all rest on treating data as vectors in a (possibly enormous)
vector space and finding a good \emph{basis} -- the very next ideas in this book. The
reason a single piece of mathematics powers images, audio, and text is that all three are
vector spaces.
\end{applicationbox}

\begin{connectionbox}
\textbf{Differential equations live in function spaces.} Consider all solutions of
$y'' + y = 0$. If $y_1$ and $y_2$ both solve it, so does $c_1 y_1 + c_2 y_2$, because
differentiation is linear:
\[
(c_1 y_1 + c_2 y_2)'' + (c_1 y_1 + c_2 y_2)
= c_1\underbrace{(y_1'' + y_1)}_{0} + c_2\underbrace{(y_2'' + y_2)}_{0} = 0.
\]
So the solution set is closed under addition and scaling -- it is a \emph{subspace} of the
space of all smooth functions. The two solutions $\cos t$ and $\sin t$ turn out to span
it, which is exactly why ``the general solution is $c_1\cos t + c_2\sin t$.'' That phrase
from your differential-equations course is a statement about a $2$-dimensional vector
space. We will make ``basis'' and ``dimension'' precise soon, and you will see this
example again.
\end{connectionbox}

\section{Subspaces: vector spaces hiding inside vector spaces}

Most of the time we do not invent a vector space from nothing. We already have a big one
-- usually $\R^n$ -- and we care about some special collection of vectors sitting
\emph{inside} it: the solutions of a system, the outputs of a matrix, a plane through the
origin. We want to ask: is that smaller collection itself a vector space, using the same
add and scale? If so, we call it a \emph{subspace}, and we inherit every theorem for free.

\begin{intuitionbox}
Picture $\R^3$, all of $3$D space. A subspace is a perfectly flat sheet (a plane) or a
perfectly straight line that \emph{passes through the origin}, extending forever in both
directions, with no edges. The defining feature: if you stand on it and walk in any
direction allowed by the sheet, or scale your position toward or past the origin, you stay
on the sheet. The origin must be included, because scaling any vector by $0$ lands you
there.
\end{intuitionbox}

\subsection{The subspace test}

We do not want to re-verify all ten axioms every time. The good news: most of them are
\emph{inherited} automatically from the big space. Commutativity, associativity, and the
distributive laws are already true for the parent's vectors, so they are true for any
subset of them. The only things that can go wrong are closure and the presence of $\vec{0}$
and negatives. That collapses to a short, practical test.

\begin{formalbox}
\textbf{Subspace test.} A nonempty subset $W$ of a vector space $V$ is a \emph{subspace}
if and only if:
\begin{enumerate}[label=(\alph*)]
  \item $\vec{0} \in W$ \hfill (it contains the origin);
  \item $\vec{u},\vec{v}\in W \implies \vec{u}+\vec{v}\in W$ \hfill (closed under addition);
  \item $\vec{v}\in W,\ c\in\R \implies c\,\vec{v}\in W$ \hfill (closed under scaling).
\end{enumerate}
Equivalently and most compactly: $W$ is a subspace iff it is nonempty and for all
$\vec{u},\vec{v}\in W$ and all scalars $c,d$, the linear combination
$c\,\vec{u} + d\,\vec{v}$ is again in $W$. ``Closed under linear combinations'' is the
whole condition.
\end{formalbox}

Why is condition (a) really needed if we have (b) and (c)? Because (c) with $c = 0$ forces
$\vec{0}\in W$ \emph{provided $W$ has at least one vector}. So (a) is really just the
reminder ``$W$ is not empty,'' and the fast way to disqualify a candidate is to check
whether $\vec{0}$ is in it. If the origin is missing, stop -- it is not a subspace.

\begin{trapbox}
\textbf{The number-one mistake: forgetting the origin.} A line that does \emph{not} pass
through the origin -- say all points with $x_2 = x_1 + 1$ in $\R^2$ -- is not a subspace,
even though it is perfectly straight and ``looks linear.'' The point
$\vec{0}=\colvec{0\\0}$ is not on it ($0 \ne 0 + 1$), and adding two of its points lands
you off the line. Always, \emph{always} check the origin first. ``Through the origin'' is
the dividing line between a subspace and a mere shifted copy (an \emph{affine} set, which
we meet when solving $\mat{A}\vec{x}=\vec{b}$ with $\vec{b}\ne\vec{0}$).
\end{trapbox}

\subsection{A worked subspace and a worked non-subspace}

\begin{example}\label{ch11:line-is-subspace}
\textbf{A line through the origin is a subspace.} Let
\[
W = \left\{ \colvec{t\\2t} : t\in\R \right\} \subseteq \R^2,
\]
the line $x_2 = 2x_1$. Check the three conditions.
\begin{itemize}
  \item \emph{Origin:} take $t=0$ to get $\colvec{0\\0}\in W$. Good.
  \item \emph{Addition:} $\colvec{s\\2s} + \colvec{t\\2t} = \colvec{s+t\\2(s+t)}$, which
  has the form $\colvec{r\\2r}$ with $r = s+t$. Still in $W$.
  \item \emph{Scaling:} $c\colvec{t\\2t} = \colvec{ct\\2(ct)}$, again of the form
  $\colvec{r\\2r}$. Still in $W$.
\end{itemize}
All three hold, so $W$ is a subspace -- a one-dimensional subspace, the line of slope $2$.
\end{example}

\begin{example}\label{ch11:not-subspace}
\textbf{The first quadrant is not a subspace.} Let
$W = \left\{\colvec{x_1\\x_2} : x_1\ge 0,\ x_2\ge 0\right\}$, the corner of the plane where
both coordinates are nonnegative. The origin is in it, and the sum of two nonnegative
vectors is nonnegative, so (a) and (b) hold. But scaling fails: take
$\vec{v}=\colvec{1\\1}\in W$ and $c=-1$. Then $c\vec{v} = \colvec{-1\\-1}$, which has
negative entries and is \emph{not} in $W$. One counterexample to one condition is enough.
The first quadrant is not a subspace.
\end{example}

\begin{figure}[H]
\centering
\begin{tikzpicture}[scale=0.92,>=Latex]
  % LEFT: line through origin = subspace
  \begin{scope}[shift={(-4.6,0)}]
    \draw[help lines,gray!25] (-2.2,-2.2) grid (2.2,2.2);
    \draw[->] (-2.4,0)--(2.4,0); \draw[->] (0,-2.4)--(0,2.4);
    \draw[very thick,targetgreen] (-1.1,-2.2)--(1.1,2.2);
    \fill[targetgreen] (0,0) circle (1.6pt);
    \draw[->,vecblue,thick] (0,0)--(0.8,1.6);
    \draw[->,vecred,thick] (0,0)--(-0.8,-1.6);
    \node[targetgreen] at (0,-2.9) {\small subspace};
    \node at (0,-3.4) {\small line through $\vec{0}$};
  \end{scope}
  % MIDDLE: line not through origin = NOT subspace
  \begin{scope}[shift={(0,0)}]
    \draw[help lines,gray!25] (-2.2,-2.2) grid (2.2,2.2);
    \draw[->] (-2.4,0)--(2.4,0); \draw[->] (0,-2.4)--(0,2.4);
    \draw[very thick,trapred] (-2.2,-1.2)--(1.2,2.2);
    \fill[white,draw=trapred,thick] (0,0) circle (1.8pt);
    \node[trapred] at (0,-2.9) {\small NOT a subspace};
    \node at (0,-3.4) {\small misses $\vec{0}$};
  \end{scope}
  % RIGHT: quadrant = NOT subspace
  \begin{scope}[shift={(4.6,0)}]
    \draw[help lines,gray!25] (-2.2,-2.2) grid (2.2,2.2);
    \fill[trapred!18] (0,0)--(2.2,0)--(2.2,2.2)--(0,2.2)--cycle;
    \draw[->] (-2.4,0)--(2.4,0); \draw[->] (0,-2.4)--(0,2.4);
    \draw[->,vecblue,thick] (0,0)--(1,1);
    \draw[->,trapred,thick,dashed] (0,0)--(-1,-1);
    \node[trapred] at (0,-2.9) {\small NOT a subspace};
    \node at (0,-3.4) {\small scaling escapes};
  \end{scope}
\end{tikzpicture}
\caption{Left: a line through the origin is closed under add and scale -- a subspace.
Middle: a parallel line missing the origin is not. Right: the first quadrant contains the
origin and is closed under addition, but scaling by $-1$ escapes it (dashed arrow). All
three conditions must hold.}
\label{ch11:fig-subspace-zoo}
\end{figure}

\begin{walkbox}
Take a walk and classify these subsets of $\R^2$ in your head -- subspace or not, and
which condition fails if not: (i) the single point $\{\vec{0}\}$; (ii) all of $\R^2$;
(iii) the $x$-axis; (iv) the unit circle $x_1^2 + x_2^2 = 1$; (v) the two coordinate axes
together (their union); (vi) all vectors with $x_1 + x_2 = 0$. Two of these are sneaky.
The circle misses the origin \emph{and} fails closure; the union of the axes contains the
origin and is closed under scaling but \emph{not} under addition -- add $\colvec{1\\0}$ and
$\colvec{0\\1}$ and you fall off both axes.
\end{walkbox}

\section{The complete list of subspaces of $\R^3$}

A beautiful thing happens in $\R^3$: there are only four \emph{kinds} of subspace, and you
can picture every one of them. This is a great reality check on the whole idea.

\begin{theorem}\label{ch11:r3-subspaces}
Every subspace of $\R^3$ is exactly one of the following:
\begin{enumerate}[label=(\roman*)]
  \item the \textbf{zero subspace} $\{\vec{0}\}$ -- just the origin (dimension $0$);
  \item a \textbf{line through the origin} (dimension $1$);
  \item a \textbf{plane through the origin} (dimension $2$);
  \item \textbf{all of $\R^3$} (dimension $3$).
\end{enumerate}
There is nothing else. No curves, no half-spaces, no off-center planes.
\end{theorem}

We are not ready to prove this rigorously -- it needs the idea of dimension, a few chapters
ahead -- but the picture is already convincing. A subspace must contain whole lines
through the origin (closure under scaling) and the parallelogram-sums of any two such lines
(closure under addition). The moment a subspace contains two vectors pointing in genuinely
different directions, it must contain the entire flat plane they span. The moment it
contains three vectors in genuinely different directions, it is forced to be all of $\R^3$.
The ``size'' can only jump $0,1,2,3$.

\begin{figure}[H]
\centering
\begin{tikzpicture}[scale=1.5,>=Latex,
    line/.style={very thick}]
  % axes
  \draw[->,gray] (0,0,0)--(1.6,0,0) node[right]{\small$x$};
  \draw[->,gray] (0,0,0)--(0,1.6,0) node[above]{\small$z$};
  \draw[->,gray] (0,0,0)--(0,0,1.8) node[below left]{\small$y$};
  % plane through origin (a parallelogram patch)
  \filldraw[targetgreen!22,draw=targetgreen,thick]
     (-1.1,-0.55,0.9)--(1.1,0.55,0.9)--(1.1,0.55,-0.9)--(-1.1,-0.55,-0.9)--cycle;
  \node[targetgreen] at (1.25,0.8,-0.9) {\small plane};
  % line through origin
  \draw[line,structurepurple] (-1.0,-1.0,0.3)--(1.0,1.0,-0.3);
  \node[structurepurple] at (1.15,1.05,-0.3) {\small line};
  % origin
  \fill[vecred] (0,0,0) circle (1.0pt);
  \node[vecred] at (-0.18,-0.18,0) {\small$\vec{0}$};
\end{tikzpicture}
\caption{The nontrivial subspaces of $\R^3$: lines and planes through the origin. Every
subspace contains $\vec{0}$ (red), and every line and plane shown extends infinitely and
is perfectly flat. Adding the zero subspace $\{\vec{0}\}$ and all of $\R^3$ completes the
list.}
\label{ch11:fig-r3subs}
\end{figure}

\begin{example}\label{ch11:plane-subspace}
\textbf{A plane through the origin, checked.} Let
$W = \left\{ \colvec{x\\y\\z} : x + 2y - z = 0 \right\}\subseteq\R^3$. This is the plane
with normal vector $\colvec{1\\2\\-1}$ through the origin.
\begin{itemize}
  \item \emph{Origin:} $0 + 2(0) - 0 = 0$, so $\vec{0}\in W$.
  \item \emph{Addition:} if $x_1+2y_1-z_1=0$ and $x_2+2y_2-z_2=0$, then summing,
  $(x_1+x_2)+2(y_1+y_2)-(z_1+z_2)=0$, so the sum satisfies the equation.
  \item \emph{Scaling:} multiplying the equation by $c$ gives $cx+2(cy)-(cz)=0$, so
  $c\vec{v}\in W$.
\end{itemize}
A subspace. Contrast with the \emph{affine} plane $x+2y-z = 5$: the origin fails
($0\ne 5$), and adding two solutions gives a right-hand side of $10$, not $5$. Not a
subspace. The pattern is general: a single homogeneous equation
($=0$, no constant term) through the origin gives a subspace; a nonzero constant on the
right breaks it.
\end{example}

\begin{focusbox}
The phrase to burn in: \textbf{homogeneous means through the origin means subspace.} The
solution set of $\mat{A}\vec{x}=\vec{0}$ is always a subspace (we will call it the
\emph{nullspace}). The solution set of $\mat{A}\vec{x}=\vec{b}$ with $\vec{b}\ne\vec{0}$ is
\emph{never} a subspace -- it is a shifted copy of one. This single distinction organizes
the next several chapters.
\end{focusbox}

\section{Spans are subspaces}

We close with the construction that makes subspaces appear everywhere. Recall the
\emph{span} of a set of vectors is all of their linear combinations:
\[
\spn\{\vec{v}_1,\dots,\vec{v}_k\} = \{ c_1\vec{v}_1 + \cdots + c_k\vec{v}_k : c_1,\dots,c_k\in\R \}.
\]
We claimed earlier, when we first met span, that it produces lines and planes through the
origin. Now we can say exactly why: \emph{a span is automatically a subspace.} This is the
single most common way subspaces arise.

\begin{theorem}\label{ch11:span-is-subspace}
For any vectors $\vec{v}_1,\dots,\vec{v}_k$ in a vector space $V$, the set
$\spn\{\vec{v}_1,\dots,\vec{v}_k\}$ is a subspace of $V$. In fact it is the
\emph{smallest} subspace containing all of the $\vec{v}_i$.
\end{theorem}

\begin{proof}
Write $S = \spn\{\vec{v}_1,\dots,\vec{v}_k\}$. We run the subspace test.

\emph{Origin.} Choosing all coefficients zero gives
$0\vec{v}_1 + \cdots + 0\vec{v}_k = \vec{0}\in S$.

\emph{Closed under addition.} Take two members of $S$,
\[
\vec{u} = a_1\vec{v}_1 + \cdots + a_k\vec{v}_k, \qquad
\vec{w} = b_1\vec{v}_1 + \cdots + b_k\vec{v}_k.
\]
Adding and grouping by each $\vec{v}_i$,
\[
\vec{u} + \vec{w} = (a_1+b_1)\vec{v}_1 + \cdots + (a_k+b_k)\vec{v}_k,
\]
which is again a linear combination of the $\vec{v}_i$, hence in $S$.

\emph{Closed under scaling.} For any scalar $c$,
\[
c\,\vec{u} = (ca_1)\vec{v}_1 + \cdots + (ca_k)\vec{v}_k \in S.
\]

All three conditions hold, so $S$ is a subspace. For ``smallest'': any subspace $W$ that
contains all the $\vec{v}_i$ must (by closure under linear combinations) contain every
combination $c_1\vec{v}_1+\cdots+c_k\vec{v}_k$, i.e. all of $S$. So $S\subseteq W$ for every
such $W$; $S$ is the smallest.
\end{proof}

\begin{intuitionbox}
The proof is almost too easy, and that is the point. ``A linear combination of linear
combinations is again a linear combination.'' Adding two recipes that each only use the
ingredients $\vec{v}_1,\dots,\vec{v}_k$ produces a recipe that still only uses those
ingredients. You can never escape the span by combining things inside it -- which is
precisely closure.
\end{intuitionbox}

\begin{example}\label{ch11:span-plane}
The span of $\vec{v}_1 = \colvec{1\\0\\1}$ and $\vec{v}_2 = \colvec{0\\1\\1}$ in $\R^3$ is
the set of all $\colvec{a\\b\\a+b}$. By the theorem it is a subspace, and since the two
vectors point in different directions it is a plane through the origin
(Figure~\ref{ch11:fig-span-plane}). Which plane? A
vector $\colvec{x\\y\\z}$ lies on it exactly when $z = x + y$, that is
$x + y - z = 0$ -- the homogeneous plane equation with normal $\colvec{1\\1\\-1}$. The
span description and the equation description are two faces of the same subspace;
converting between them is what elimination is for, coming up next.
\end{example}

\begin{figure}[H]
\centering
\begin{tikzpicture}[scale=1.4,>=Latex]
  % axes
  \draw[->,gray] (0,0,0)--(1.7,0,0) node[right]{\small$x$};
  \draw[->,gray] (0,0,0)--(0,1.7,0) node[above]{\small$z$};
  \draw[->,gray] (0,0,0)--(0,0,1.9) node[below left]{\small$y$};
  % the spanned plane z = x + y, drawn as the parallelogram of combos a*v1 + b*v2.
  % corners are combinations with a,b in {-1.2, 1.2}; tikz order is (x, z, y).
  % a*v1+b*v2 = (a, b, a+b) in math (x,y,z) -> tikz (x=a, z=a+b, y=b)
  \filldraw[targetgreen!20,draw=targetgreen,thick]
     (1.2,2.4,1.2)--(-1.2,0,1.2)--(-1.2,-2.4,-1.2)--(1.2,0,-1.2)--cycle;
  % v1 = (1,0,1): x=1,z=1,y=0  -> tikz 3d order is (x,z,y)
  \draw[->,very thick,vecblue] (0,0,0)--(1,1,0) node[right]{\small$\vec{v}_1$};
  % v2 = (0,1,1): x=0,z=1,y=1
  \draw[->,very thick,vecred] (0,0,0)--(0,1,1) node[above right]{\small$\vec{v}_2$};
  % a sample combination v1+v2 = (1,2,1)
  \draw[->,thick,structurepurple,dashed] (0,0,0)--(1,2,1);
  \node[structurepurple] at (1.15,2.05,1) {\small$\vec{v}_1+\vec{v}_2$};
  \fill[black] (0,0,0) circle (0.9pt);
  \node at (-0.16,-0.2,0) {\small$\vec{0}$};
\end{tikzpicture}
\caption{The span of two independent vectors $\vec{v}_1,\vec{v}_2$ in $\R^3$ is the whole
plane through the origin that they lie in -- here the plane $x+y-z=0$. Every linear
combination (such as the dashed $\vec{v}_1+\vec{v}_2$) stays in the green sheet: that is
Theorem~\ref{ch11:span-is-subspace}, drawn.}
\label{ch11:fig-span-plane}
\end{figure}

\begin{labbox}
By hand, with $\vec{v}_1=\colvec{2\\1}$ and $\vec{v}_2=\colvec{-1\\3}$ in $\R^2$:
\begin{enumerate}
  \item Compute $2\vec{v}_1 + 3\vec{v}_2$ and confirm it is a vector in $\R^2$ (closure).
  \item Argue that since $\vec{v}_1,\vec{v}_2$ point in different directions, their span is
  all of $\R^2$. (Hint: can you hit $\colvec{1\\0}$? Solve
  $c_1\vec{v}_1+c_2\vec{v}_2=\colvec{1\\0}$.)
  \item Now replace $\vec{v}_2$ by $\colvec{-4\\-2}$. What is $\spn\{\vec{v}_1,\vec{v}_2\}$
  now, and what dimension is it? (The second vector is $-2\vec{v}_1$, so it adds nothing.)
\end{enumerate}
This little lab is the seed of \emph{independence} and \emph{dimension}: whether the
spanning vectors genuinely add new directions.
\end{labbox}

\begin{strangbox}
When Strang introduces the column space $\col(\mat{A})$ as ``all combinations of the
columns of $\mat{A}$,'' he is taking a span -- so by Theorem~\ref{ch11:span-is-subspace}
it is automatically a subspace, no further checking needed. Likewise the nullspace is a
span of the special solutions. Half of Strang's ``four subspaces'' chapter is really this
one theorem applied four times. Carry it with you.
\end{strangbox}

\section{Operations on subspaces}

Two subspaces can be combined to make new ones, and one combination that looks natural
secretly fails. This sharpens your instinct for closure.

\begin{proposition}\label{ch11:intersection}
If $U$ and $W$ are subspaces of $V$, then their intersection $U\cap W$ (vectors in both)
is a subspace.
\end{proposition}

\begin{proof}
The origin lies in both $U$ and $W$, so $\vec{0}\in U\cap W$. If
$\vec{u},\vec{v}\in U\cap W$, then both lie in $U$ (so $\vec{u}+\vec{v}\in U$ and
$c\vec{u}\in U$) and both lie in $W$ (so $\vec{u}+\vec{v}\in W$ and $c\vec{u}\in W$). Hence
the sum and the scaling lie in the intersection. Subspace test passed.
\end{proof}

Geometrically, two planes through the origin in $\R^3$ meet in a line through the origin
(or coincide); a plane and a line meet in the origin or the whole line. Every such
intersection is again a flat object through the origin -- a subspace, as promised
(Figure~\ref{ch11:fig-intersection}).

\begin{figure}[H]
\centering
\begin{tikzpicture}[scale=1.3,>=Latex]
  \draw[->,gray] (0,0,0)--(1.7,0,0) node[right]{\small$x$};
  \draw[->,gray] (0,0,0)--(0,1.6,0) node[above]{\small$z$};
  \draw[->,gray] (0,0,0)--(0,0,1.9) node[below left]{\small$y$};
  % first plane (sourceblue patch)
  \filldraw[sourceblue!18,draw=sourceblue,thick]
     (-1.3,-0.7,1.1)--(1.3,0.7,1.1)--(1.3,0.7,-1.1)--(-1.3,-0.7,-1.1)--cycle;
  % second plane (targetgreen patch), tilted the other way
  \filldraw[targetgreen!18,draw=targetgreen,thick]
     (-1.1,1.0,1.1)--(1.1,-1.0,1.1)--(1.1,-1.0,-1.1)--(-1.1,1.0,-1.1)--cycle;
  % their intersection: a line through the origin
  \draw[very thick,structurepurple] (0,0,-1.6)--(0,0,1.7);
  \node[structurepurple] at (0.05,0.28,1.7) {\small$U\cap W$};
  \fill[black] (0,0,0) circle (0.9pt);
  \node at (0.18,-0.22,0) {\small$\vec{0}$};
\end{tikzpicture}
\caption{Two distinct planes through the origin in $\R^3$ (blue and green) intersect in a
line through the origin (purple). The intersection of subspaces is always a subspace
(Proposition~\ref{ch11:intersection}); here it is forced to be one-dimensional.}
\label{ch11:fig-intersection}
\end{figure}

\begin{trapbox}
\textbf{Unions usually fail.} The \emph{union} $U\cup W$ (vectors in either one) is almost
never a subspace. Take the $x$-axis and the $y$-axis in $\R^2$: both are subspaces, but
their union is not closed under addition, since
$\colvec{1\\0}+\colvec{0\\1}=\colvec{1\\1}$ lies on neither axis. The fix is to take not
the union but the \emph{sum} $U + W = \{\vec{u}+\vec{w} : \vec{u}\in U,\ \vec{w}\in W\}$,
which \emph{is} a subspace (and for the two axes equals all of $\R^2$). Intersection: yes.
Sum: yes. Plain union: no.
\end{trapbox}

\section{Looking ahead}

You now have the core vocabulary. A \emph{vector space} is any world closed under
adding and scaling; a \emph{subspace} is a smaller such world living inside a bigger one,
detected by the three-line test with the origin as your first checkpoint. \emph{Spans are
always subspaces}, which is how subspaces actually show up -- as the column space, the
nullspace, and every line and plane through the origin.

The questions that drive the next chapters are now askable in clean language. Given a
subspace, how few vectors do we need to span it (a \emph{basis})? How many is that, exactly
(its \emph{dimension})? When does a list of vectors waste effort by repeating directions
(\emph{dependence})? And how do the four subspaces attached to a matrix $\mat{A}$ fit
together? Every one of those is a question about closure, and you now know what closure
means.

\begin{aestheticbox}
Step back and notice what just happened. We defined ``space'' without ever saying what a
vector \emph{is}. A vector became anything that adds and scales lawfully -- an arrow, a
matrix, a polynomial, a sound, an image, a solution of a differential equation. By
refusing to look inside the vectors, we made our theorems apply to all of them at once.
That refusal -- proving things from rules rather than from pictures of coordinates -- is
the move that turns linear algebra from arithmetic into a theory, and it is the same move
that lets the same eigenvalue you compute for a $2\times 2$ matrix tell you the resonant
frequency of a bridge, the steady state of a web-ranking algorithm, and the principal
direction of a cloud of data.
\end{aestheticbox}

\begin{recoverybox}
\textbf{Intuition recap -- carry these five lines.}
\begin{itemize}
  \item A \emph{vector space} is any world closed under \emph{add} and \emph{scale}; the
  eight bookkeeping axioms come free, closure is the real test.
  \item A \emph{subspace} is a vector space living inside another one. Test it in three
  lines: contains $\vec{0}$, closed under $+$, closed under scaling -- equivalently, closed
  under linear combinations.
  \item First reflex on any candidate: \emph{is the origin in it?} If not, stop -- not a
  subspace. ``Homogeneous ($=0$) means through the origin means subspace.''
  \item A \emph{span} is automatically a subspace (Theorem~\ref{ch11:span-is-subspace}), in
  fact the smallest one containing the given vectors. This is how column spaces and
  nullspaces are born.
  \item \emph{Intersection} of subspaces: yes. \emph{Sum}: yes. Plain \emph{union}:
  almost never.
\end{itemize}
There is no arithmetic ``cost'' to tally here -- the subspace test is the cheapest check in
the book, three closure questions -- but the conceptual payoff is the largest: prove it
once for vector spaces, and it holds for arrows, matrices, polynomials, functions, and data.
\end{recoverybox}

\section*{Exercises}
\addcontentsline{toc}{section}{Exercises}

\begin{exercise}\label{ch11:ex-classify-r2}
For each subset of $\R^2$, decide whether it is a subspace. If not, name a specific
condition that fails with a numerical counterexample.
(a) $\{\colvec{x\\y} : y = 3x\}$; \quad (b) $\{\colvec{x\\y} : y = 3x + 1\}$;
(c) $\{\colvec{x\\y} : xy = 0\}$; \quad (d) $\{\colvec{x\\y} : x = y = 0\}$;
(e) $\{\colvec{x\\y} : x \ge 0\}$.
\end{exercise}

\begin{exercise}\label{ch11:ex-r3-plane}
Show that $W = \{\colvec{x\\y\\z} : 2x - y + 3z = 0\}$ is a subspace of $\R^3$ by checking
the three conditions, and find two specific vectors whose span is $W$. What is the normal
vector to this plane?
\end{exercise}

\begin{exercise}\label{ch11:ex-degree-two}
Explain carefully why the set of polynomials of degree \emph{exactly} $2$ is not a
subspace of $P_2$. Give two failures: one about the zero vector, and one about closure
under addition (find two degree-$2$ polynomials whose sum has degree less than $2$).
\end{exercise}

\begin{exercise}\label{ch11:ex-matrix-subspace}
Inside $M_{2\times 2}$, consider the set of all matrices with trace zero, i.e.
$\begin{bmatrix}a&b\\c&d\end{bmatrix}$ with $a + d = 0$. Show it is a subspace. Then do the
same for the symmetric matrices ($b = c$). For each, find a spanning set.
\end{exercise}

\begin{exercise}\label{ch11:ex-span-equal}
In $\R^3$, let $\vec{v}_1 = \colvec{1\\1\\0}$, $\vec{v}_2 = \colvec{0\\1\\1}$,
$\vec{v}_3 = \colvec{1\\2\\1}$. Show that $\vec{v}_3 \in \spn\{\vec{v}_1,\vec{v}_2\}$ and
conclude that $\spn\{\vec{v}_1,\vec{v}_2,\vec{v}_3\} = \spn\{\vec{v}_1,\vec{v}_2\}$. What
kind of subspace of $\R^3$ is this (point, line, plane, or all)?
\end{exercise}

\begin{exercise}\label{ch11:ex-intersection}
Let $U$ be the $xy$-plane and $W$ the plane $x = z$ in $\R^3$. Describe $U \cap W$
explicitly (find a vector that spans it) and confirm geometrically that it is a line
through the origin. Why does Proposition~\ref{ch11:intersection} guarantee in advance that
the intersection is a subspace?
\end{exercise}

\begin{exercise}\label{ch11:ex-functions}
Let $V$ be the space of all functions $f:\R\to\R$. For each subset, decide whether it is a
subspace and justify with the closure conditions:
(a) all $f$ with $f(0) = 0$; \quad (b) all $f$ with $f(0) = 1$;
(c) all $f$ with $f(1) = f(2)$; \quad (d) all even functions ($f(-x)=f(x)$);
(e) all $f$ with $\int_0^1 f(x)\,dx = 0$.
\end{exercise}

\begin{exercise}\label{ch11:ex-union-proof}
Prove the general fact behind the union trap: if $U$ and $W$ are subspaces of $V$ and
\emph{neither} contains the other, then $U \cup W$ is not a subspace.
(Hint: pick $\vec{u}\in U\setminus W$ and $\vec{w}\in W\setminus U$ and consider
$\vec{u}+\vec{w}$. Show it cannot lie in $U$ and cannot lie in $W$.)
\end{exercise}

\begin{exercise}\label{ch11:ex-axiom-derive}
Using only the ten axioms, prove that the zero vector is unique: if $\vec{0}$ and
$\vec{0}'$ both satisfy $\vec{v}+\vec{0}=\vec{v}$ and $\vec{v}+\vec{0}'=\vec{v}$ for all
$\vec{v}$, then $\vec{0}=\vec{0}'$. (Hint: evaluate $\vec{0}+\vec{0}'$ two ways.) Then
prove negatives are unique.
\end{exercise}

\begin{exercise}\label{ch11:ex-smallest}
Theorem~\ref{ch11:span-is-subspace} says $\spn\{\vec{v}_1,\dots,\vec{v}_k\}$ is the
\emph{smallest} subspace containing the $\vec{v}_i$. Use this to give a one-line argument
that the intersection of \emph{all} subspaces containing a set $S$ equals $\spn(S)$. (This
is a slick alternative definition of span: ``the smallest subspace containing $S$.'')
\end{exercise}

\begin{exercise}\label{ch11:ex-de-subspace}
\textbf{(Connection to calculus.)} Show that the set of all differentiable functions $y$
satisfying $y' = 3y$ is a subspace of the space of differentiable functions. Then find a
nonzero function in it. (You may recall the solution is $y = Ce^{3t}$ -- which says the
subspace is the span of a single function. What is its dimension?)
\end{exercise}
